\documentclass[11pt]{article}
\usepackage[margin=1in]{geometry}
\usepackage{enumitem}
\usepackage{booktabs}
\usepackage{hyperref}
\setlist{itemsep=2pt,topsep=3pt}

\title{Independent Review and Revision Plan for ACM TACO\\
\large Shape Sets the Chunk: Selecting and Keying the Chunk of Chunked
Linear-Recurrence GPU Kernels}
\date{Cycle taco-r2, opened 2026-10-16}

\begin{document}
\maketitle
\sloppy

\section*{Scope and Inputs}

This file opens revision cycle taco-r2. It synthesizes a fresh, independent
five-reviewer panel (Editor-in-Chief, R1 methodology and statistics, R2
domain, R3 practitioner perspective, devil's advocate) and a fresh re-check of
every taco-r1 item. The reviewers read the current manuscript (ACM acmsmall,
main article 25 pages of which 20 precede the references; Supplementary
Material 31 pages) jointly, from the sources, the compiled PDFs and the result
files only, without reference to earlier review rounds, against the standard
of ACM Transactions on Architecture and Code Optimization (new regular paper;
at most 20 pages excluding references; double-anonymous review). The full
reports are in \texttt{docs/review\_reports/taco\_r2\_*.md}. Each concern below
cites the findings it merges: EIC M$n$/m$n$, R1 M$n$/m$n$, R2 M$n$/m$n$, R3
M$n$/m$n$ (major/minor issue $n$ of that report), DA C$n$ (devil's-advocate
finding), ED (editorial decision) and RC (re-check verdict on a taco-r1 item).
Line numbers refer to the current \texttt{paper/main.tex} and
\texttt{paper/appendix.tex}.

Effort classes: (a) text only; (b) offline analysis of existing result files;
(c) new GPU measurement, with an estimate in GPU-hours. Class (c) is split by
host in this cycle:
\begin{itemize}
  \item \textbf{(c, RTX 3090 lab host)}: runs on the lab host with four stock
        RTX~3090 GPUs (Ampere, sm\_86). These boards differ from the modified
        48\,GB RTX~3090 of the existing data, so their results are reported as
        a separate, stock platform (with driver, library versions, clocks,
        power limit, host CPU and whether access is exclusive) and, where they
        repeat an earlier measurement, as a cross-instance check. Hardware
        counters are collected there only where the profiler is permitted.
  \item \textbf{(c, optional: needs a rented data-centre GPU)}: needs a
        serving-engine deployment or a data-centre GPU (H100/H800/H200/B200).
        The author chose not to rent one in this cycle, so these items are
        optional; if one is not run, the claim it would support is scoped as
        the item states, in the abstract, the conclusion and Limitations.
\end{itemize}
Items of class (c) are run only after the author approves the host and the
budget.

\section*{Editorial Decision}

\textbf{Major Revision.} Four reviewers recommend Major Revision (EIC, R1, R2,
R3). The devil's advocate recommends Reject, with resubmission of a refocused
paper possible, on two CRITICAL findings (DA C1, low significance of a
confirmatory result; DA C2, the headline chunk sensitivity belongs to the
in-house kernel), and would move to Major Revision given engine-level evidence
and a demoted in-house kernel.

\begin{center}
\small
\begin{tabular}{lrrrrrl}
\toprule
Reviewer & Novelty & Method. & Exper. & Clarity & Signif.
  & Recommendation \\
\midrule
EIC & 45 & 74 & 68 & 38 & 45 & Major Revision \\
R1 Methodology & 55 & 60 & 66 & 48 & 52 & Major Revision \\
R2 Domain & 35 & 65 & 60 & 40 & 40 & Major Revision \\
R3 Practitioner & 52 & 72 & 63 & 38 & 45 & Major Revision \\
Devil's advocate & 30 & 60 & 55 & 35 & 28 & Reject (refocus) \\
\midrule
\textbf{Mean} & \textbf{43.4} & \textbf{66.2} & \textbf{62.4} & \textbf{39.8}
  & \textbf{42.0} & 4 $\times$ Major, 1 $\times$ Reject \\
\bottomrule
\end{tabular}
\end{center}

Against the taco-r1 panel (means 42.4 / 60.4 / 61.2 / 40.0 / 40.4) methodology
improved by about 6 points and the other dimensions are flat; no dimension
regressed by more than 0.2 points.

\paragraph{Re-check of taco-r1.} Of 66 items, 47 are addressed, 11 partial, 8
declined with scoped claims and 0 not addressed. Acceptance criteria: 5 of 8
met, 3 partial (A1 artifact anonymity; A5 latency-weighted a-priori headline;
A8 supplement length and notation). Stage 11 gate: page limit, PDF and source
anonymity, references and citation keys pass; the artifact archive fails
anonymity. The partial items are folded into the concerns below.

\paragraph{Consensus strengths (keep).} Recurrence-preserving timing with fp32
output checks; held-out prompt-cluster evaluation with bootstrap, TOST and
permutation tests; the placement-versus-values and warm/cold-L2 interventions;
the compiler-specialization bit in the key; candid reporting of model
failures; the autotuner-state and output-divergence controls. All of about 70
recomputed numbers matched the result files, with two definitional
differences (93.5\% vs.\ 94\% tied cells; 0.40 vs.\ 0.38\% between two tables).

\paragraph{Consensus weaknesses.} (i) Significance and novelty: for
data-oblivious kernels a shape-only optimum is the null expectation, the lemma
is elementary, the refuted design is one unreviewed preprint, and the positive
production-kernel result is small (1.0--5.9\% of scan time in isolated
CUDA-graph replay) with wide intervals for its recoverable part; (ii) the steep
chunk effect belongs to the in-house Triton kernel, whose launch table couples
the chunk with grid size and warps and whose ``default 256'' is borrowed;
(iii) no engine-level or data-centre evidence for the device-bound claims;
(iv) the cost model does not pay on the paper's own numbers and its a-priori
accuracy depends on the selection protocol; (v) clarity (mean 39.8): five or
more timing conventions, several definitions of ``default'' and ``best
static'', two figures in 20 pages, about 50 pointers into a 31-page
supplement; (vi) measurement reliability: a launch-inclusive session that did
not reproduce, unguarded data sets, a possibly virtualized and power-capped
RTX~3090, untested SSD/Gated DeltaNet value invariance above the floor, and
H800 input effects absent from the main text.

\paragraph{What separates Major from Minor Revision.} The text and offline
items address clarity, fair baselines and most methodological points but are
unlikely on their own to lift significance above the bar. The decisive lever
is evidence that the device-bound chunk effect matters end to end, or a clear
demonstration that it does not, reframed as the finding (M1d, M3b, M3e).

% ======================================================================
\section*{Venue Compliance (MUST FIX)}

% STAGE: 11 (integrity gate) + 10 (revision)
% RULES: academic-paper/agents/formatter_agent.md
\begin{enumerate}[label=V\arabic*.]
  \item \textbf{Artifact anonymity (Stage 11 gate FAIL).} The availability
        statement (main.tex:1455--1461) promises an anonymized archive that
        does not yet exist in a clean, current form. In the raw result-file
        folder, 122 files contain absolute user paths, 148 git hashes, 146
        host names, 166 the project name and 6 a user name, and the directory
        name is the previous venue's; together these name the system of the
        cited preprint and a previous venue. The latest Code Ocean capsule
        (built 2026-10-10) has removed user names but still carries
        previous-venue directory names and a real host name in two files, and
        it predates the taco-r1 runs, so it does not back Table~7 or the
        supplementary packed-call, autotuner-state and divergence tables.
        Rebuild the capsule from the current result files with host names,
        paths, hashes, project and venue names and package version suffixes
        stripped, rename the directories, include every file behind a main or
        supplementary table, and re-run an anonymization scan before upload.
        (a)+artifact rebuild. (EIC M8; RC V2; ED)
  \item \textbf{Paper-folder hygiene.} The paper folder holds author-named
        profile photographs (the \texttt{figs/profile} subfolder), backups for
        three prior venues and a previous venue's checklist (the
        \texttt{backup} and \texttt{legacy} subfolders), unused old figures
        from the earlier framing, and figure PDFs whose creation dates carry a
        local timezone offset (not propagated into main.pdf). Upload only the
        two PDFs for review; if sources are ever required, build a clean
        bundle with main.tex, appendix.tex, ref.bib and the referenced figures
        only. (a). (RC V2 and Stage 11 gate; ED)
  \item \textbf{Confidential disclosure of the preprint relationship.} The
        refuted design is a public preprint (arXiv:2604.10597), cited in the
        third person, which ACM permits. If it is the authors' own work, say
        so to the editor in the confidential cover letter; confirm that the
        footnote ``New Paper, Not an Extension of a Conference Paper''
        (main.tex:91) remains accurate; and make sure any public code
        repository linked from the preprint does not expose this manuscript or
        its result files during review. (a). (EIC M8; R2 M6 and question 5;
        RC further observations)
  \item \textbf{Page limit kept.} The limit currently passes: Discussion,
        Limitations, Conclusion and the availability and AI-use disclosure end
        on page 20, and ``References'' is the first heading of page 21. Every
        addition requested below (claims-to-evidence table, mechanism figure,
        deployment-effect table, promoted supplementary results) must be
        offset by the cuts of M1b and M8b; keep the AI-use disclosure, the
        availability statement, CCS concepts, keywords, running head and the
        new-paper footnote inside the 20 pages. (a). (RC Stage 11 gate)
  \item \textbf{Supplement length and dependence.} The supplement is 31 pages
        with more than 40 tables, against the taco-r1 target of about 20; the
        main text points into it about 50 times. Reduce it to about 20--25
        pages: merge the convention, launch-inclusive oracle and interval
        tables (Tables S on conventions, oracle\_li, tables\_ci, ci), drop
        secondary statistics that no main-text claim needs, and keep what
        reproduction requires. (a). (RC M9c, A8; EIC M5; R3 M7; ED)
\end{enumerate}

% ======================================================================
\section*{Major Concerns}

\subsection*{M1. Significance and framing (CRITICAL, DA C1)}
Merges EIC M3, DA C1, DA C8, DA C16, R2 M6, R1 m4, EIC m9, EIC m16, ED reason
1. All five reviewers: the three kernels contain no data-dependent branches,
early exits or sparse paths (main.tex:386--391), so a shape-only optimum is
the null expectation; the lemma is one inequality; the only foil is one
unreviewed preprint whose map selects chunk 256 for essentially every call
(Fig.~S6); and shape-keyed caches (Triton autotune, Kernel Launcher, FFTW
wisdom, cuDNN benchmark mode, TensorRT timing caches) already embody the
recipe. The findings that are new to TACO readers are buried.

% STAGE: 10 (revision)
% RULES: academic-paper/agents/draft_writer_agent.md
\begin{enumerate}[label=M1\alph*.]
  \item (a) \textbf{Contributions = the robust findings.} Make the
        contributions (i) the key must be the kernel's compiler
        specialization signature (one divisibility bit changes the compiled
        kernel and its optimal chunk; a length-bucket key without it leaves
        22--90\,$\mu$s of headroom); (ii) the execution regime (host-bound
        vs.\ device-bound) decides whether the chunk matters; (iii) buffer
        placement and a warm L2 masquerade as data or chunk effects in tuning
        (median 21\%, up to 43.6\%, against 0.1\% for values); (iv) models
        calibrated in eager timing fail under CUDA graphs; (v) the economics
        of profiling versus modelling. Add a short ``what was not known
        before'' paragraph that positions each finding against TFLA, the
        Triton autotuner, Kernel Launcher, input-aware autotuning (Nitro, Ding
        et al., Seer) and the production practice of M9b. \emph{Where:} title
        and abstract main.tex:39--70; Introduction 112--139; Contributions
        141--162; Sec.~2 264--280, 322--327. (EIC M3; DA C1; R2 M6; ED)
  \item (a) \textbf{Condense the refutation to one subsection.} State Lemma
        3.1 as an observation or remark (main.tex:364--404); keep the
        placement/L2 finding and one predictor table; move the routed study
        (1177--1194), the break-even and key-churn derivation (1291--1318) and
        the information-theoretic detail of Sec.~6.3 (876--902) to the
        supplement; keep the split-call timing pitfall (523--532) as a
        two-sentence protocol note that does not read as a correction of the
        cited preprint. (R2 M6; DA C16; EIC M3; EIC m16)
  \item (a)+(b) \textbf{Headroom as an information result with a host-side
        cost corollary.} ``Below one host synchronization at every tuned
        shape'' holds for the exact-shape table on 29 covered shapes; the
        recommended bucket key leaves 26\,$\mu$s on one RTX~3090 key, and
        bucket 2048, batch 4, unaligned leaves 30--32\,$\mu$s, above one
        synchronization (21.4\,$\mu$s); under real traffic the exact-length
        key puts 51.5--83.8\% of calls on keys below the break-even (Tables
        S37, S38). The device-side reduction cost is quoted for the RTX~3090
        only (7.3\,$\mu$s at $L=2048$, 42\,$\mu$s at $L=8192$); on the RTX
        PRO~6000 it is $+2.7$ and $-3.7\,\mu$s, zero within noise, so ``$g>o$
        is out of reach'' cannot be claimed there for device-side policies.
        Report the RTX PRO~6000 cost; support ``every tuned shape'' with
        per-shape upper confidence bounds (reference scan: held-out per-shape
        maximum 14.5 against 21\,$\mu$s); say that an isolated synchronization
        is a lower bound on its critical-path cost in device-bound pipelines;
        separate the single-static-chunk baseline of main.tex:1211--1216 from
        the shape-keyed baseline of the lemma's $h(s)$; restate the abstract
        and Sec.~6.7 accordingly. (DA C8; R1 m4; EIC m9)
  \item (a) \textbf{Author decision: the significance claim.} Fix the paper's
        central claim after M3b. If a shape-keyed chunk table changes a
        graph-captured prefill by about 1\% or less end to end, make that the
        finding (for production kernels the chunk is a second-order knob end
        to end; keying and measurement pitfalls are what matter) and shorten
        the paper accordingly, as R2 and DA C1 route (b) suggest; if the
        effect is larger, lead the abstract with it. \emph{Where:} abstract;
        Conclusion main.tex:1445--1453. \emph{Requires an author decision.}
        (ED; R2 recommendation; DA C1)
\end{enumerate}

\subsection*{M2. The in-house kernel and its baselines (CRITICAL, DA C2)}
Merges EIC M2, R1 M7, R2 M1, R3 M6, DA C2, C3, C11, C19, RC (Shapley
convention), ED reason 2. The steep chunk curve, the 49\% per-call default
regret (34\% of scan time), the 8.5$\times$ spread, the 4.4--25\,ms per
forward pass of Table~8, the tensor-parallel width sensitivity and the whole
cost model belong to the Triton Mamba-1 scan written for this study. It runs
at 0.81--1.54$\times$ the reference CUDA scan (geometric mean 1.23, slower at
10 of 12 keys; Table S9, RTX~3090 only); its launch table was chosen once on a
fourth GPU (an RTX~3070) and used unchanged (S1.2); the best table entry is
7--18\% behind full launch tuning (up to 53\% at one key); chunks 512 and 1024
sit at 255 registers with 12 and 64\,B of spill (Table S13); its ``default
256'' is borrowed from Mamba-2. Every independently written kernel shows the
default or a stock length rule within a few percent: reference scan 3.1\%
(device time, Table S21) and 6.6\% (launch-inclusive, Table~6), TFLA 1.4\%
(Table S27), SSD and Gated DeltaNet host-bound at most 2.5\%; above the
device-time floor SSD chunk 256 is the oracle for 375 of 397 calls and Gated
DeltaNet chunk 64 for 154 of 155.

% STAGE: 10 (revision) + src
% RULES: academic-paper/agents/argument_builder_agent.md
\begin{enumerate}[label=M2\alph*.]
  \item (a) \textbf{Disclose provenance and role.} State in the abstract and
        Introduction that the Triton scan is the authors' own, 1.23$\times$
        the reference scan, with a launch table chosen on another GPU; present
        it in every caption as a steep-curve stress case; rename it (e.g.\
        ``Triton scan (ours)'' vs.\ ``CUDA scan (mamba\_ssm)'') instead of
        redefining ``Mamba-1'' for the rest of the paper (main.tex:653--660).
        Correct ``register-bound'': at batch 1 (65\% of the calls) the kernel
        launches only $\lceil d/D_b\rceil\times B$ programs and runs 2--6 warps
        per SM, so it is grid- and latency-bound there. \emph{Where:}
        abstract; Sec.~4 (406--518); Sec.~5.1 (536--545); Sec.~6 opening.
        (DA C2; EIC M2; R2 M1.2, M7; R3 M6; R3 m3)
  \item (b) \textbf{Fair untuned baselines.} Use as untuned references the
        free length rule $\mathrm{clip}(\mathrm{nextpow2}(L),32,1024)$, which
        mirrors the stock CUDA dispatch (device time, per call 5.86 / 7.76 /
        10.66\%, latency-weighted 5.96 / 16.31 / 13.94\% on the RTX PRO~6000 /
        RTX~3090 / H800), the best chunk per batch (20.0 / 17.4 / 24.7\%) and
        a length-bucket table (2--20\,$\mu$s per call, Table S23). Restate the
        Sec.~6 opening, the abstract and Tables 4, 5 and 8 against them and
        keep chunk 256 (48.6 / 96.4 / 50.9\%) as a secondary row; the best
        static chunk by latency on the RTX PRO~6000 (512) loses 15.5\% of scan
        time, not 34\%. Against the free rule the calibrated model buys only
        1.7 / 1.5 / 7.7 points per call; explain why the calibration-only
        length tree (16.6 / 20.4 / 38.4\%, Table S19) is far worse than this
        free rule. For the reference scan, use its stock length rule as the
        Mamba-1 default. \emph{Where:} main.tex:653--660; Tables 4, 5, 8
        (785--859, 1242--1269). (DA C3; EIC M2; R2 M1.5; R3 M6; ED)
  \item (a) \textbf{One definition of ``default'' and ``best static''.}
        Table~7's ``Default'' is the library default for SSD and Gated
        DeltaNet but the best single chunk for the Triton scan; Table~8's best
        static is chosen by total latency per batch, Table~5's held out by
        per-call regret; the Sec.~6 opening mixes both (``20\% (27\%)'').
        Define each once per kernel, label every exception, and state up
        front that the production-kernel tuning gains are those of Table~7.
        (DA C19, C12; R3 M6.3; EIC M2)
  \item (b)+(c, RTX 3090 lab host) \textbf{Chunk effect at fixed launch
        configuration.} The default table maps chunks 32/64/128 to $(D_b,
        n_w)=(8,2)$, 256/512 to $(4,2)$ and 1024 to $(4,4)$, so from chunk 128
        to 256 the grid doubles (256 to 512 programs at $d=2048$, batch 1).
        From the existing 184-configuration sweep at the twelve keys, report
        per GPU the spread over chunks within each fixed $(D_b, n_w,
        \mbox{stages})$, the movement of the optimum and the best-chunk
        regret at fixed configuration. R2's offline check of the short search
        timings: on the RTX~3090 at fixed $(4,2)$ chunk 128 is fastest at 11
        of 12 keys and one chunk loses about 0\%, against an optimum over
        $\{32,64,128,512\}$ and 17\% along the default table; on the H800 at
        fixed $(8,2)$ chunk 128 wins at 11 of 12 keys (1.2\%) against 21\%;
        on the RTX PRO~6000 at fixed $(4,2)$ the optimum still moves.
        Re-time that slice under the full protocol on the lab host (about 2--3
        GPU-hours). Define the ``29 + 3 points'' decomposition (a Shapley
        split; 32.7 = 29.3 + 3.4) and give it in device time, the primary
        convention (RTX PRO~6000: 50.6 points chunk, $-6.2$ launch). State in
        Table~1 and Sec.~6.2 that chunk results are results for launch-table
        entries. \emph{Where:} Sec.~4 (421--423); main.tex:968--977; Table~1;
        Fig.~S1; Tables S7, S13. (R1 M7; R2 M1.1, fix (a); R3 M6.4; RC
        Shapley convention)
  \item (a)+(b) \textbf{Chunk-only scope.} The best full configuration differs
        between warm and cold L2 at 7 of 12 keys, and the cost model used as a
        full launch-space tuner loses 8.9--23\% on average (53--95\% maximum)
        when one key is left out (Table S17); production caches tune the full
        configuration per key. Either justify the chunk-only scope (library
        APIs such as the SSD and FLA entry points expose only the chunk) and
        restrict the claims, or report chunk selection jointly with full
        configuration tuning from the existing sweeps. \emph{Where:}
        Sec.~6.4 ``Auto-tuning baseline''; Recipe; Limitations main.tex:1425.
        (DA C11; R2 M1 fix (b))
  \item (c, optional, RTX 3090 lab host) \textbf{Parity or per-GPU retuning.}
        Either retune the launch table per GPU before studying the chunk, or
        bring the kernel closer to the reference (a slice-based carry instead
        of the masked full-tile reduction
        \texttt{h\_last = sum(where(t == C-1, hs, 0))} and the per-state
        masked reduction; more programs per channel block at small $B\cdot d$;
        an aligned main loop with a masked tail), and show whether the chunk
        sensitivity persists. About 3--6 GPU-hours plus kernel engineering.
        If not done, M2a's stress-case framing is mandatory. (R2 M1.3--1.4,
        fix (c); EIC M2 fix (c); DA C2)
\end{enumerate}

\subsection*{M3. Production relevance and device-bound evidence}
Merges EIC M1, R1 M1, R2 M4, R2 M5, R3 M1, R3 M2, R3 M3, DA C5, C6, C10, RC
M6a, M6e, ED reason 3. The only positive production-kernel result is that
under isolated CUDA-graph replay the defaults lose 10--24\% per call but
1.0--5.9\% of scan time (RTX PRO~6000: 1.75 / 5.85 / 3.19\%; RTX~3090: 2.80 /
2.20 / 1.02\% for SSD 370M / SSD 2.7B / Gated DeltaNet 1.4B), i.e.\
0.10--1.44\,ms per forward pass, measured with one graph per (call, chunk) on
two non-data-centre GPUs, without the scan's share of the forward pass and
without an engine. R3's estimate from projection GEMMs alone puts the
default's cost at about 0.3\% (2.7B) to 0.7\% (370M) of an average forward
pass; 99 / 92 / 100\% of the default's absolute regret (SSD 2.7B / SSD 370M /
Gated DeltaNet 1.4B) comes from calls under 4096 tokens.

% STAGE: 10 (revision) + src
% RULES: academic-paper/agents/argument_builder_agent.md
\begin{enumerate}[label=M3\alph*.]
  \item (b) \textbf{Recoverable gains with intervals.} The 1.0--5.9\% is the
        default's latency-weighted excess over a noisy per-call oracle. Report
        instead the recoverable share (default minus held-out shape table)
        with 42-cluster bootstrap intervals in Table~7, the abstract and the
        conclusion. R1's recomputation (graph replay, warm L2, three rounds):
        \begin{center}\small
        \begin{tabular}{lrrl}
        \toprule
        Kernel, GPU & Default LW & Held-out table LW & Recoverable [95\% CI] \\
        \midrule
        SSD 370M, RTX 3090 & 2.80\% & 0.01\% & 2.79\% [1.51, 5.30] \\
        SSD 370M, RTX PRO 6000 & 1.75\% & 0.03\% & 1.72\% [0.89, 3.57] \\
        SSD 2.7B, RTX 3090 & 2.20\% & 0.03\% & 2.17\% [1.15, 4.61] \\
        SSD 2.7B, RTX PRO 6000 & 5.85\% & 0.46\% & 5.39\% [2.77, 7.63] \\
        GDN 1.4B, RTX 3090 & 1.02\% & 0.57\% & 0.45\% [$-0.90$, 1.78] \\
        GDN 1.4B, RTX PRO 6000 & 3.19\% & 0.25\% & 2.94\% [1.88, 4.93] \\
        \bottomrule
        \end{tabular}
        \end{center}
        The lower end of the abstract's range (Gated DeltaNet on the RTX~3090)
        is not a gain a table delivers, and ``recovers to within 0.7\%''
        refers to per-call means, not scan time. Scope the conclusion to the
        two GPUs measured (main.tex:1447--1449) and add ``no whole-model graph
        or serving-engine measurement'' to Limitations if M3b and M3e are not
        run. \emph{Where:} abstract main.tex:59--62; Contributions; Sec.~6.5,
        Table~7 and Fig.~2 (1096--1150); Conclusion. (R1 M1; RC M6a, M6e; ED)
  \item (c, RTX 3090 lab host) \textbf{Full-model prefill as one CUDA graph.}
        Capture a full-model prefill forward pass at fixed shapes (Mamba-2
        370M and 2.7B, Gated DeltaNet 1.4B) as one CUDA graph, once with the
        library default chunk and once with the held-out shape-keyed table;
        measure prefill latency in at least two sessions; report the measured
        scan share of the forward pass and a main-text ``deployment effect''
        table (oracle scan time, default and table regret in ms, share of the
        full forward pass) at three or four representative tokens-per-call
        points, next to ``\% of scan time''. About 4--8 GPU-hours. This is the
        minimum end-to-end evidence; it decides M1d. \emph{Where:} Table~7;
        Sec.~6.5; Limitations main.tex:1425--1429. (R1 M1 fix (ii); EIC M1 fix
        (c); R3 M1 fix 1; R2 M4.6; ED)
  \item (a) \textbf{Engine execution modes and regime names.} Sec.~1 states
        without citation that ``production engines run prefill both host-bound
        and as captured CUDA graphs''. In vLLM V1's default
        \texttt{FULL\_AND\_PIECEWISE} mode, prefill and mixed batches are not
        captured as full graphs, the Mamba and linear-attention mixers are
        splitting ops that stay eager, and Mamba backends list full-graph
        support only for uniform single-token decode. Describe, per engine and
        mode, which prefill scan calls run device-bound and at which token
        counts (graph capture, C++ launch paths, host run-ahead under async
        scheduling), and qualify or drop the general claim. The host-bound
        regime is defined by the Python launch path (56--90\,$\mu$s per Triton
        call; an SSD floor of about 0.46--0.5\,ms on the RTX PRO~6000 and
        0.74--1.5\,ms on the RTX~3090, against 33--57\,$\mu$s under graph
        replay); rename it, e.g.\ ``eager, launch-bound (Python launch
        path)'', and report the harness's per-call launch overhead next to
        it. Mention speculative-decoding verification and chunked-prefill
        slices as sources of very short calls. \emph{Where:} Sec.~1;
        Conventions box main.tex:608--642; Sec.~6.5; Recipe. (R2 M4.1--4.2;
        R3 M2; DA C6; R3 m23)
  \item (c, RTX 3090 lab host) \textbf{Packed steps as engines run them.} The
        one serving-like experiment contradicts the recipe: on SSD 370M a
        held-out packed-token-bucket table is worse than the default (fresh
        state: 9.75\% per step, 7.46\% latency-weighted, against 14.14\% and
        10.23\%; adding the number of sequences gives 13.46\%); with carried
        state 10.68\% (8.11\%) falls only to 9.21\% (6.77\%); SSD 2.7B falls
        from 9.04 to 2.77\% (fresh) and 7.53 to 4.48\% (carried). The 370M
        data are noise-dominated: the per-round argmin is identical in all
        three rounds for 34\% of SSD 370M step variants (70\% for 2.7B, 83\%
        for Gated DeltaNet), the median round-to-round deviation is 1.4\%
        (90th percentile 6.3\%), and the steps were timed in device time
        behind the 20\,ms spin that S1.1 and S8 reject for multi-kernel
        pipelines. Re-time the packed steps under CUDA-graph replay with at
        least five rounds and a noise floor (split-half or a null as in S12);
        capture packed variable-length steps at padded capture sizes (16 to
        512 tokens and up to the largest prefill budget an engine would
        capture), replay them with varying composition, and report the regret
        of a table keyed on capture size alone and of richer keys
        (per-sequence length structure). Apply the carried state the same way
        for SSD and Gated DeltaNet (now the first sequence vs.\ every
        sequence) or justify the difference. Revise Contribution~2 and the
        recipe's packed-length key to what the data support. About 3--5
        GPU-hours. \emph{Where:} Sec.~6.5 ``Packed variable-length calls''
        (1152--1162); Recipe (1320--1344); S10, Table S35. (DA C5; R3 M2 fix
        2; R3 m20; R2 m8)
  \item (c, optional: needs a rented data-centre GPU) \textbf{Serving-engine
        integration.} Integrate a shape-keyed chunk table into vLLM or SGLang
        with the engine's default CUDA-graph configuration, for one Mamba-2
        hybrid (e.g.\ Nemotron-H 8B, Falcon-H1 or Granite-4.0-H) and one Gated
        DeltaNet hybrid (a small Qwen3-Next-style model or the 1.4B
        checkpoint), with a realistic chunked-prefill token budget on public
        traces; report time to first token, throughput, the share of mixer
        calls that are device-bound and below 4096 tokens, and the switching
        overhead (metadata rebuild, graph count); include one TP~$=2$
        configuration. About 15--25 GPU-hours. If not run, the abstract and
        conclusion must say that no engine-level measurement was made and that
        the device-bound results are isolated-kernel replays. (R2 M4 fix; DA
        C6; EIC M1 fix (b); ED)
  \item (c, optional: needs a rented data-centre GPU)
        \textbf{Data-centre replication.} Repeat Table~7 (graph replay of SSD
        370M/2.7B and Gated DeltaNet 1.4B), the packed-call study and the
        autotuner-state control on a stock data-centre GPU, recording the
        board power limit and locking application clocks where permitted.
        About 6--10 GPU-hours. If not run: state in Limitations that no
        data-centre-class (HBM) GPU was replayed under graphs, that Triton
        code generation on sm\_120 and sm\_86 may differ from sm\_90/sm\_100
        for these matrix-product-heavy pipelines, and scope every
        production-kernel claim to ``an Ampere consumer GPU and a Blackwell
        server GPU''. (EIC M1 fix (a); R2 M5; R3 M3.1; DA C10; RC M6a)
  \item (c, optional, RTX 3090 lab host) \textbf{Production kernels at shard
        layouts.} Tensor parallelism is measured only for the Mamba-1
        kernels. Time SSD at $n_{\mathrm{heads}}/n_{\mathrm{tp}}$ with
        $\mathrm{ngroups}>1$ (as in Nemotron-H or Falcon-H1) and Gated
        DeltaNet at Qwen3-Next-like head counts divided by $n_{\mathrm{tp}}$;
        the four GPUs of the lab host also allow a TP~$=2$--4 run with
        collectives active. About 2--4 GPU-hours. (R2 M4.5; DA C10)
\end{enumerate}

\subsection*{M4. The cost model}
Merges EIC M4, R1 M3, R2 M3, R3 M5, DA C4, RC M3c, M4b, M5a, M2a residuals,
N8, N9, ED reason 4. On the paper's own numbers the model does not pay:
calibration costs 130 / 322 / 202\,s with a warm compile cache (493 / 843 /
664\,s cold) against 3.9 / 7.4 / 6.1\,s to profile every chunk of all 57 exact
keys; compilation dominates (363 of 493\,s), so the model saves only about
0.03--0.05\,s of timing per key and pays only beyond 1622--2475 keys; the
serving replays produce 11--23 bucket keys per trace, and non-enumerable keys
cover at most about 2\% of calls under bucketed keys (Table S38). Its
a-priori accuracy depends on the selection protocol, simple rankers and MLPs
tie or win on some GPU, and transfer to another kernel is untested.

% STAGE: 10 (revision) + src
% RULES: academic-paper/agents/argument_builder_agent.md
\begin{enumerate}[label=M4\alph*.]
  \item (b) \textbf{One pre-specified a-priori protocol; all six reported.}
        The released analysis evaluates six protocols (candidate sets
        ablation12 and ablation12+work\_only, times the criteria LOGO on the
        other GPUs' captured calls, LOGO on calibration LOSO, and
        target-calibration LOSO). In device time they give: RTX PRO~6000
        2.24, 0.38, 0.56, 4.16, 0.38, 0.56\%; RTX~3090 6.23, 2.49, 1.74, 6.23,
        2.49, 0.88\%; H800 2.14--3.01\%; overall 0.38--6.2\%. The paper
        reports one (ablation12+work\_only, LOGO captured calls), the most
        pessimistic on the RTX PRO~6000; the target-calibration LOSO
        criterion, which is the procedure the ``What transfers'' paragraph
        recommends, gives 0.56 / 0.88 / 2.14\%, and the tie rule alone moves
        the RTX PRO~6000 between 0.38, 0.56 and 4.16\%. Pre-specify and
        justify one protocol; report all six in a supplementary table with
        latency-weighted regret, maxima (53 / 27 / 50\% for the reported
        forms), per-batch values and cluster intervals; list the candidate
        forms and their scores per fold. Replace ``3.0 to 6.2\%'' in the
        abstract and conclusion by the pre-specified protocol with an interval
        or by the protocol range, each with its latency-weighted companion: on
        the RTX PRO~6000 the a-priori form (the three-parameter restriction)
        loses 26.9\% of scan time latency-weighted (CI 23.5--29.3), no better
        than the best chunk per batch (27.0\%); the frozen form loses 2.3 /
        9.7 / 5.3\% (RTX~3090: 9.7\% latency-weighted against 3.5\% per call).
        Report per-batch regrets (RTX~3090: 15.9\% at batch 16 against 1.5\%
        at batches 1 and 4; H800: 6.3\% at batch 4) and aligned/unaligned
        regrets (0.22/1.01, 3.89/2.25, 2.69/4.24\%). \emph{Where:} abstract
        main.tex:64--65; Sec.~6.4 (907--960, 923--930); Table~5; Tables
        S14--S19; Conclusion 1450--1451. (R1 M3; DA C4; RC M4b, N8 and
        advisory; R2 m10)
  \item (b)+(a) \textbf{Fair comparison of calibration-only selectors.}
        Evaluate every calibration-only selector under the same a-priori
        protocol in the main table. Under it the eight-parameter model has the
        worst GPU mean (4.16 / 6.23 / 3.01\%, mean 4.47, worst 6.23) against
        its three-parameter restriction (4.20 / 0.85 / 3.53, mean 2.86), the
        MLP over five seeds (1.59 / 1.31 / 7.13, mean 3.34) and the structural
        MLP (3.75 / 4.90 / 1.11, mean 3.25); ``within 3.5\% on every GPU'' and
        ``best on the RTX PRO~6000'' use the frozen form selected on two
        GPUs' calls. State plainly that on the only form-blind GPU (H800) the
        full form is not better than the restriction (paired $-0.53$ [$-1.35$,
        $+0.60$], Table S15) and is beaten by the structural MLP (0.57\%; 1.1\%
        over five seeds); across GPUs it is better only on the RTX PRO~6000
        ($-3.80$ [$-4.78$, $-2.89$]) and worse on the RTX~3090 ($+2.66$
        [$+2.12$, $+3.17$]). The H800 nine-parameter variant (add\_l2\_miss)
        reproduces the frozen form exactly (3.0077\%), so it is not
        independent evidence. Move the three-parameter top-3 result into
        Table~5 so that the abstract claim has a main-text anchor. (DA C4; R1
        M3 fix (iii); RC M9c)
  \item (b) \textbf{Fit documentation.} Report: the refit on full-protocol
        timings alone (0.73 / 4.28 / 6.05\%; the H800 doubles from 3.0\%);
        inverse-variance weighting (fit unchanged); the objective $R^2$ (0.89 /
        0.93 / 0.94); the largest search-versus-full-protocol discrepancy
        (226\%, launch-inclusive RTX~3090); that 2208 of 2298 calibration
        timings are brief 3/10 search timings taken in key order without
        randomization (drift confounded with key) and, on two GPUs, without
        the guard; the 12-key bootstrap with the RTX PRO~6000 point estimate at
        the lower edge of [0.38, 2.26] (RTX~3090 [1.74, 6.09]); the
        leave-one-key-out regrets (0.8 / 7.4 / 3.7\%) as the more honest
        uncertainty with 12 clusters; and the recommended optimizer protocol,
        given that 50 random starts reach a 16\% lower objective and 1.7\%
        regret on the RTX~3090. Report calibration on the six default-table
        configurations only (the 72 static re-measurements alone reach 0.56 /
        1.8 / 6.1\%). \emph{Where:} main.tex:932--935; appendix.tex:846--848,
        858--880. (RC M3c; R1 M3; R2 M3(d); R3 m22)
  \item (a)+(b) \textbf{Physical consistency of Eq.~(3).} The latency term
        $A\cdot\min(1, W_1/W)$ lets the predicted wave time per tile drop below
        $A$, the serial per-tile latency of the non-unrolled state loop, so
        the model can predict wave makespans below the critical path; the
        fitted $W_1$ is 6.6--10.5 warps, exceeded in batched calls with 4-warp
        configurations. Justify the form or refit a
        $\max(\mbox{latency},\mbox{throughput})$ form (Hong and Kim; Volkov)
        and report whether the selected chunks change. Model the last partial
        wave ($m=\min(r,\lceil n_p/n_{\mathrm{SM}}\rceil)$ assumes an even
        spread; $n_p/n_{\mathrm{SM}}$ is 1.4--2.7 at batch 1); validate the
        static spill-size proxy $\omega$ against measured local-memory traffic
        at the 255-register configurations ($C=512, 1024$; see M5e); and
        correct appendix.tex:793, which still calls $\sigma$ ``spill-store
        traffic''. \emph{Where:} Sec.~4, Eq.~(3) (406--518). (R2 M3(b), m5,
        m6; RC M3b)
  \item (a)+(b) \textbf{Reposition the model.} Either give a concrete case
        where keys cannot be enumerated, or shorten Sec.~4, Sec.~6.4 and the
        abstract to a compact case study whose conclusion is that profiling
        dominates for enumerable keys. The one use profiling cannot replace is
        an ahead-of-time table for a GPU that cannot be timed: move the S2
        cross-device results for all source/target pairs into Sec.~6.4 (RTX
        PRO~6000 parameters with H800 compile facts select with 3.8\%, 3.3\%
        in device time). Present the model as a top-$k$ pruner and compare it
        with Triton's perf\_model/top\_k mechanism and an off-the-shelf
        occupancy-only ranker. Drop the Gated DeltaNet ``instance'' from
        Table~2 (five parameters fitted to nine timings, no regret). Soften
        appendix.tex:742 (``a new GPU needs its compile facts but hardly its
        own calibration'', one kernel on three GPUs) and main.tex:1449--1451
        (add ``for this kernel''). (R3 M5; EIC M4 fix (b); R2 M3(c); DA C4; RC
        M2a residuals)
  \item (b)+(c, RTX 3090 lab host) \textbf{Work-term model for SSD and Gated
        DeltaNet.} Add the padding term for $L<C$ ($\max(L,C)$ work, as
        disclosed at main.tex:1424) and report eager and graph regrets; the
        eager-calibrated model loses 9--22\% under graphs and the replay-fold
        refit loses 20\% of scan time on SSD 370M on the RTX PRO~6000 against
        1.7\% for the default. Calibrate it on graph-timed calibration keys
        (about 1 GPU-hour on the lab host) instead of replay folds.
        \emph{Where:} Table~7 work-term columns; S9. (RC N9, M6a; DA C4; R2
        M2)
  \item (c, optional, RTX 3090 lab host) \textbf{Transfer to a second
        kernel.} Apply the unchanged procedure (compile without running, read
        the resource report, calibrate at about twelve keys, terms fixed
        before seeing the target calls) to an independently written Triton
        kernel whose registers and spills vary with the chunk, e.g.\ FLA's
        \texttt{chunk\_fwd\_o} or the SSD \texttt{chunk\_scan} kernel with the
        inner tile fixed. About 4--8 GPU-hours. If not run, keep the
        single-kernel case-study scoping (main.tex:406--414, 501--506).
        (R2 M3(e); EIC M4 fix (a); DA C4)
  \item (c, optional: needs a rented data-centre GPU) \textbf{A fourth,
        unseen GPU.} Freeze all selectors and the protocol and evaluate on a
        previously unseen GPU (A100, L40S or H100 SXM); this is the only way to
        support ``a new GPU needs its compile facts but hardly its own
        calibration''. About 2--4 GPU-hours. If not run, remove that sentence.
        (R1 M3 fix (iv))
\end{enumerate}

\subsection*{M5. Measurement reliability and statistics}
Merges R1 M2, M4, M5, EIC M6, R2 M2, DA C7, C10, RC M4e, M6d, M10c, ED reason
6, R1 m5--m7, m11, m13, EIC m12, m13, R2 m3.

% STAGE: 10 (revision) + src
% RULES: academic-paper/agents/argument_builder_agent.md
\begin{enumerate}[label=M5\alph*.]
  \item (b) \textbf{Noise floors and floor sensitivity.} Next to every regret,
        report the noise floor of that data set (cross-round or cross-session
        regret of the per-shape best chunk, as already computed for the TP
        sweep and in Table S37) and interpret effects only above it; say in
        each caption that regrets against the per-call oracle are biased up by
        the minimum of noisy means, and that the held-out tables' 0.06--0.5\%
        in device time and graph mode bounds this bias. Give a sensitivity
        analysis (factors 1.5 / 2 / 3) for the ad hoc launch floor (``twice the
        median batch-1 latency at $L\le256$'') and the 100\,$\mu$s convention
        threshold, which drive counts such as 375 of 397 calls and 17\% of
        calls above the floor. For ``no practically relevant information'',
        count only data sets whose chunk effect exceeds the noise floor.
        \emph{Where:} Tables 4, 5, 7, 8; main.tex:1088--1092; Table S37. (R1
        M2 fix (i), m5, m6; EIC m12; RC N2)
  \item (a)+(b) \textbf{The RTX~3090 tensor-parallel sweep is
        noise-dominated.} A table retuned at the exact shard shape and scored
        on another round loses 14.87\% on the Triton kernel, and the cost
        model beats it by $-8.35$ points [$-12.38$, $-4.59$] (reference scan
        $-4.35$ [$-7.06$, $-1.86$]), which only scoring noise can produce; the
        per-round argmin agrees for 52\% of shapes. On the H800 the retuned
        table loses 0.75\% and the model is worse by $+2.16$ [1.16, 3.22], as
        expected. Demote the RTX~3090 TP numbers (including ``a table tuned at
        the unsharded width loses 36.8\%'') to the supplement as
        uninformative and rest the shard-width claim on the H800, or re-run
        the sweep with more rounds and the guard on the lab host (c, optional,
        RTX 3090 lab host; about 2--3 GPU-hours). \emph{Where:} Sec.~6.5
        ``Tensor-parallel widths'' (1062--1077); Contribution~2; Tables S19,
        S24. (R1 M2; DA C10; R3 M3)
  \item (c, RTX 3090 lab host) \textbf{Value invariance above the floor.}
        Table~3 (414 cells) is launch-inclusive, unguarded and at
        $L=512$ launch-dominated (median 62\,$\mu$s at the selected chunk,
        inside the 56--90\,$\mu$s launch cost); the captured-activation input
        was timed in a second process 36 minutes after the others, which
        confounds it with placement; with three rounds only 4 of 18
        configurations pass TOST at a 2\% margin, the main text reports median
        bounds (3.6 / 3.2\%) while the maxima reach 14.8\% (Mamba-1) and
        34.0\% (SSD) on the RTX PRO~6000; at $L=2048$ the values-only ranges
        reach 22.8\% but are non-significant through noise, and the 0.1\% vs.\
        21\% headline is made at $L=8192$ only. SSD and Gated DeltaNet
        invariance above their floor is untested (two of three device-time SSD
        data sets are inconclusive at the 0.5-point margin), and the Gated
        DeltaNet synthetic sweep calls a different entry point than the
        captured calls. Re-run in device time or graph replay with all eight
        inputs including captured activations interleaved in one process,
        fixed placement (design A) plus the multi-allocation design C, at
        least five rounds, Mamba-1 at $L\in\{2048, 8192\}$, and SSD and Gated
        DeltaNet above their floor through the same entry point for synthetic
        and captured inputs; report per-configuration and worst-case TOST
        bounds in the main text. About 3--4 GPU-hours. Until done, qualify the
        recipe's ``synthetic inputs suffice'' for SSD and Gated DeltaNet
        (M7c). \emph{Where:} Sec.~6.1, Table~3 (695--744); Tables S1, S5,
        S10; Limitations. (R1 M5; DA C7; ED; RC recipe-vs-Limitations)
  \item (a) \textbf{H800 input effects in the main text.} On the H800, 28 of
        36 fixed-chunk configurations show a significant input effect
        launch-inclusive and 36 of 36 in device time, with ranges up to 47.6\%
        and 39.9\% at $L=2048$; the per-input best chunk flips between 256 and
        1024 there, only 54--71\% of cells are tied, and the regret reaches
        5.7\%. The replication paragraph (main.tex:1386--1388) cites only the
        RTX~3090 (``at most 4.3\%''). Add one sentence: on the H800 these
        effects could not be attributed to placement because no placement or
        L2 control was run there. A placement and cold-L2 control on the H800
        is (c, optional: needs a rented data-centre GPU; about 1--2
        GPU-hours). \emph{Where:} main.tex:742--744, 1386--1395;
        appendix.tex:710. (DA C7; RC M6d)
  \item (c, RTX 3090 lab host) \textbf{Hardware counters for the mechanism.}
        The binding resource per regime (grid parallelism, register-limited
        occupancy, spill traffic, issue slots, launch floor) is asserted, and
        the placement explanation rests on an intervention only. With Nsight
        Compute where permitted (cloud containers return ERR\_NVGPUCTRPERM),
        collect for the twelve tuning keys at each chunk: achieved occupancy,
        warps per scheduler, issue-slot utilization, top stall reasons,
        local-memory load/store sectors and L2 hit rate; the L2 hit rate for
        two placements of the same input; and a per-kernel time breakdown
        versus chunk under graph replay at short and long calls for SSD
        (chunk\_cumsum, chunk\_state, state\_passing, bmm\_chunk, chunk\_scan)
        and Gated DeltaNet (cumsum, kkt, solve\_tril, recompute\_w\_u, fwd\_h,
        fwd\_o), including the empty or masked programs when $C>L$. These feed
        the mechanism figure (M8d) and the spill-proxy check (M4d). About 2--4
        GPU-hours. \emph{Where:} Sec.~4; Table~1; Sec.~6.1 ``Placement, not
        values''; Sec.~6.5. (R2 M2; ED; R2 m5)
  \item (c, RTX 3090 lab host) \textbf{Platform corroboration and session
        reproducibility.} The 48\,GB RTX~3090 is non-stock, possibly
        virtualized or time-sliced (main.tex:559--561), with the power cap
        active in 38--92\% of readings, and it alone carries the routed study,
        the reference-scan sweep, the larger checkpoints, the production
        kernels at larger widths and TFLA. Launch-inclusive results did not
        reproduce between two RTX PRO~6000 sessions (median oracle 60 vs.\
        98\,$\mu$s; shape tree 0.7 to 4.7\%; model 0.9 to 4.9\%), and the RTX
        PRO~6000 device-time sweep and graph replays are single sessions.
        State whether the 48\,GB board is a virtual GPU and which conclusions
        rest on it alone. On the lab host, re-time the RTX~3090 graph replays
        of Table~7 and a stratified subset of the captured calls and
        reference-scan sweep with exclusive access (locked clocks if
        possible), and report the cross-session stability of a tuned table
        (tune in session A, serve in session B). About 3--5 GPU-hours. A
        second RTX PRO~6000 session (device-time captured calls and graph
        replays) is optional and needs a rented RTX PRO~6000 (about one GPU
        day). \emph{Where:} Sec.~5.2 (557--606); Sec.~6.8 (1392--1395);
        S1.1, S8; Tables S1, S2. (EIC M6; R1 M2 fix (ii); R3 M3.2--3.3; DA
        C10)
  \item (a) \textbf{Protocol documentation.} Add Table S1 rows for the four
        taco-r1 runs (RTX PRO~6000 graph replays, packed calls,
        autotuner-state re-timing, output divergence) and a Table S2 entry for
        the stock 24\,GB RTX~3090 used for the last two (state they are
        cross-instance checks); flag every main-table number from an
        unguarded run (value invariance, launch-inclusive and H800 captured
        calls, most calibration data) or replace it; evaluate every row of
        Table~5 on the same data set per GPU (now three-round medians on two
        GPUs, single blocks on the H800, and the RTX~3090 a-priori row from
        another session); record host CPU model, governor and load for every
        launch-inclusive data set; document what the local mamba\_ssm build
        changes and confirm the SSD kernels are unmodified. On the 20\,ms
        device-time spin: state whether warm-ups and the enqueue of 50
        repeats of the multi-kernel pipelines (about 23--40\,ms at their eager
        floor) fall inside it, size it to the measured enqueue time in new
        runs, and compare device time with graph replay on the same Triton
        calls to bound any clock or power bias. State in the Conventions box
        that both conventions time identical kernels back to back on resident
        inputs and that cold-L2 controls exist only for Mamba-1 on the RTX
        PRO~6000 and one flushed graph round on the RTX~3090. \emph{Where:}
        Conventions box main.tex:608--642; appendix.tex:59--132, 1952--1953,
        2063. (RC M10c; EIC M6, m13; R1 M2 fix (iv)--(v), m7, m11; R2 m3)
  \item (b) \textbf{Intervals and per-batch breakdowns.} Add cluster-bootstrap
        intervals for Table~6 (reference scan, larger checkpoints) and
        Table~8; give per-batch values for the 1.4B and 2.8B transfer rows
        (the pooled 1.4B regret of 3.7\% is 1.0 / 13.9 / 0.9\% at batches 1 /
        4 / 16); report the number and size distribution of the 42 prompt
        clusters and use BCa or studentized intervals; refit inside the
        bootstrap (or repeated grouped cross-validation) for paired
        differences; control multiplicity across the about 19 data sets, 3
        conditions and 3 statistics, or label the tests descriptive.
        \emph{Where:} main.tex:1022, 1056--1058; Tables S (tables\_ci, ci)
        appendix.tex:1719--1796. (RC M4e; R1 M4 fix (iv)--(v), m13)
  \item (b)+(c, RTX 3090 lab host) \textbf{Information-side power where the
        headroom exceeds one synchronization.} On Mamba-2 370M,
        launch-inclusive RTX~3090, the oracle label is tie- and
        noise-dominated ($H(\mbox{oracle}\mid\mbox{exact key})=1.97$ bits of a
        possible 2.32; about 3.2 chunks within 2\%), so a CMI excess of
        $-0.0003$ bits ($p=0.50$) says little; R1's continuous-outcome check
        (Spearman $\rho$ $-0.06$ to 0.31, all $p\ge0.13$, 24--224 calls per
        shape) agrees but excludes only large effects; the largest headroom
        (43\,$\mu$s at $L=8192$, batch 16) is 0.4\% of a 10.7\,ms call, the
        size of the 0.9\% within-block CV. Add a continuous-outcome test
        (per-chunk log latency on entropy within the exact key with layer and
        prompt-cluster random effects, or the top-2 latency ratio) with effect
        sizes in $\mu$s, intervals and a minimum detectable effect; express
        the equivalence margin in absolute terms tied to the decision cost
        (e.g.\ half a synchronization, about 10\,$\mu$s) or relative to the
        residual regret, since the 0.5-point margin exceeds the shape tree's
        entire residual (0.38--0.46\%) and cannot fail in the beneficial
        direction. Re-time the five Mamba-2 shapes whose headroom exceeds one
        synchronization in graph mode with at least three rounds on the lab
        host (under 1 GPU-hour; a different board, so report it as
        cross-instance). \emph{Where:} Sec.~6.3; Sec.~6.7; Tables 5, S6, S37.
        (R1 M4)
\end{enumerate}

\subsection*{M6. Workload representativeness and serving economics}
Merges R1 M6, DA C9, C20, EIC M7, EIC m8, R3 M1, R3 m5, m15, m16, R2 M4.4, R1
m8, m9, R2 m14.

% STAGE: 10 (revision)
% RULES: academic-paper/agents/argument_builder_agent.md
\begin{enumerate}[label=M6\alph*.]
  \item (b) \textbf{Workload weighting.} The 2360 captured calls span only 57
        distinct (length, batch) shapes (34 / 13 / 10 lengths at batch 1 / 4 /
        16); 79.5\% have power-of-two lengths, all aligned, where real traffic
        would give about one in sixteen; 773 calls (33\%) lie on the twelve
        calibration keys; twelve shapes hold 65\% of the calls and 160 calls
        share $L=25$; the batch mix is fixed at 1536 / 488 / 336, a batch-16
        call holds 2--10 distinct prompts and 177 of 488 batch-4 calls contain
        fewer than four; batch-1 calls are 65\% of the data and 20\% of calls
        take under 20\,$\mu$s in device time on the RTX PRO~6000. Every
        per-call mean and latency-weighted share depends on this mix.
        Post-stratify the captured-call results to the (length bucket, batch)
        distributions of the Azure 2023/2024 and BurstGPT replays already
        computed in S13 and report the headline per-call and latency-weighted
        numbers under those weights, with cluster intervals, beside the
        unweighted ones. \emph{Where:} Sec.~5.2 (570--591); Tables 4, 5, 7, 8;
        abstract. (R1 M6; DA C9; EIC M7; R3 M1 fix 2)
  \item (c, optional, RTX 3090 lab host) \textbf{Natural-length calls.} Time a
        set of calls whose lengths and batch compositions are sampled from the
        public traces (natural, mostly unaligned lengths) and report the
        selectors on it. About 2--3 GPU-hours. (DA C9)
  \item (b)+(a) \textbf{Serving replay and units.} The replay costs every call,
        including packed calls, as the nearest key of the Triton scan (which
        supports neither variable-length calls nor the production Mamba-1
        path), uses a step model of 2\,ms plus 10\,$\mu$s per token that is
        2.3$\times$ faster than measured eager prefill (so the bound's
        direction depends on the baseline engine), scales $P_k$ from wall
        clocks and extrapolates beyond $L=8192$ and batch 16; the break-even
        counts (1900 / 2660 calls) and ``2.1--2.6\% / 5.5--10\% of prefill''
        are model outputs. Re-run the replay with the production kernels'
        own graph-replay and packed-step timings (170 packed steps exist on
        the RTX PRO~6000); label every measured and every modelled conclusion
        in the main text, with the load-sensitivity ranges. Define ``ms per
        forward pass'' (mean per-call regret times the number of layers over a
        non-workload-weighted mix) or switch to ms per prefill at stated
        $(L,B)$ points, and give the latency-weighted figure wherever a
        per-call regret appears in the abstract or conclusion. \emph{Where:}
        Sec.~6.7 ``Deployment units'', ``Key churn'', ``Serving traces''
        (1284--1372); Table~8 caption; appendix.tex:2206--2220. (EIC M7; R3
        M1 fix 3; R2 M4.4; R1 m9; R3 m5; R2 m14; DA C6)
  \item (a) \textbf{Routed end-to-end study.} Lead with stock dispatch against
        the best chunk ($+0.7\pm1.1$\% of prefill at 4096 tokens, Table S39),
        not with the deliberately bad chunk 128 (``adds 16.8\,ms''); write
        ``16.8\,ms (16.7\%)'' to keep units consistent; state that the testbed
        is eager Hugging Face Transformers on one RTX~3090 with the reference
        scan, that prefill takes about 46--47\,ms at 256 and 976 tokens
        whatever the chunk, and that decode at about 39--40\,ms per token
        (1.24--1.26\,s for 32 tokens) dilutes ``3.1\% of end-to-end time'';
        report prefill (time to first token) as primary. Present the
        19--191\% penalties of the statistic as an upper bound for one
        unoptimized implementation (the full histogram is not
        bandwidth-optimal, and the guarded and learned-table variants pay for
        a histogram they do not use), give a fused or device-side cost
        estimate per statistic, and note that a zero-overhead shape-keyed
        policy is never measured end to end there. \emph{Where:}
        main.tex:1177--1194; S15, Table S39. (R3 M1 fix 4, m15, m16; R1 m8;
        EIC m8; DA C20, C6)
\end{enumerate}

\subsection*{M7. Recipe correctness}
Merges R3 M4, R2 M8, R2 m11, m15, m16, DA C18, R3 m19, m21, RC
recipe-vs-Limitations. Engineers will copy the recipe box (main.tex:1320--1344);
every item must be correct and point to its evidence.

% STAGE: 10 (revision)
% RULES: academic-paper/agents/draft_writer_agent.md
\begin{enumerate}[label=M7\alph*.]
  \item (a) \textbf{Batch invariance.} ``A batch-invariant or deterministic
        mode should keep one chunk per key'' is not sufficient: the key
        contains the batch size and the packed length, so the same request
        gets a different chunk and summation order depending on what it is
        batched with, and chunk switching changes the greedy output of
        2.6--9.4\% of prompts in half precision (per-length table up to
        4.7\%; Table S36). Batch-invariant serving must fix the chunk or key
        it only on per-request quantities; replicas must share one frozen,
        versioned table. Move the divergence result into the recipe and the
        abstract-level trade-off, with its run-to-run baseline (the default
        run twice was bitwise identical). \emph{Where:} Recipe; Sec.~7
        main.tex:1410--1412; S11. (R3 M4.1; DA C18; R2 m11)
  \item (a) \textbf{Crash hazard.} Small chunks overflow the 32-bit indexing
        of the chunk-state tensor and fault with an illegal memory access (SSD
        chunk 32 at width 5120, batch 16, $L>6528$; TFLA chunks 32 and 64 at
        the largest shapes; S6, S7), which poisons the CUDA context, so a
        warm-up profiler can take a server down. Add to ``Failure modes'' a
        candidate-set rule: exclude chunks whose state tensor would exceed
        $2^{31}$ elements at the largest batch and length the key admits (the
        tensor grows as $L/C$). (R3 M4.2)
  \item (a) \textbf{Synthetic warm-up inputs.} Constant and very sparse
        inputs are systematically faster, by up to 6.7\% on the Triton scan
        and 3.4\% on the reference scan (S1.3), and placement moves latency by
        up to 43.6\%. Specify dense random inputs of realistic scale, several
        allocations or a flushed L2, and timing above the launch floor;
        qualify the sentence for SSD and Gated DeltaNet until M5c is done
        (main.tex:1328--1330 conflicts with Limitations). (R3 M4.3; RC
        recipe-vs-Limitations)
  \item (a)+(b) \textbf{The key is the specialization signature.} Triton
        specializes on the divisibility of integer arguments including
        strides and on pointer alignment, not on ``the length''; the kernel
        sees stride $=L$ only because the inputs are made contiguous, while in
        production layouts the stride is $B\cdot L$ or the packed token count,
        and $x$ comes from \texttt{xz.chunk(2, dim=1)} of a $[B,2d,L]$
        projection with an offset that may break 16-byte alignment. Define the
        key as the kernel's actual specialization signature (Triton's
        compilation cache key and specialization attributes), verify which
        attributes Triton~3.4 applies to this kernel, and report the
        aligned/unaligned effect for the SSD and Gated DeltaNet kernels on
        captured or packed calls. For the reference CUDA scan, discuss
        $\mathrm{kIsEvenLen} = (L \bmod (\mathrm{kNThreads}\cdot\mathrm{kNItems})
        = 0)$, which changes with the chunk, in Table~1 and Sec.~6.5, and
        state which part of its chunk effect is thread count, items per
        thread and this bit. Make alignment its own action: unaligned lengths
        are about 2$\times$ slower at batch 1 (Fig.~S3: $L=2382$ against
        2048), so recommend padding lengths or strides to a multiple of 16 or
        aligned allocation where the engine allows. A small aligned/unaligned
        re-timing of the production kernels is (c, RTX 3090 lab host; under 1
        GPU-hour). \emph{Where:} Sec.~3 key definition; Sec.~4 ``Compiler
        specialization and parallel shapes'' (main.tex:510--518); Table~1;
        S12. (R2 M8, m15; R3 M4.4)
  \item (a)+(c, RTX 3090 lab host) \textbf{Tuning cost under graphs.}
        ``Seconds per GPU'' holds for eager timing and excludes compilation
        (27--40\,s for the captured workload). Under graphs, tuning means one
        captured graph per candidate chunk and capture size: report the
        capture time and graph-pool memory (measure alongside M3d, under 1
        GPU-hour), or recommend device-time tuning before capture and say
        that device-time timings of multi-kernel pipelines are unreliable near
        the floor (S8). (R3 M4.6, m21)
  \item (a) \textbf{Label unsupported items as precautions.} Invalidation on a
        power-limit, driver, compiler or library change (only the GPU-model
        part is evaluated), broadcasting the chunk across TP ranks, and
        consistency with prefix-cache checkpoints have no supporting
        measurement. For the last, be concrete: engines may tie the
        state-checkpoint granularity of hybrid prefix caching and the
        per-step chunk metadata to the configured chunk, so a per-shape chunk
        may be unsupported without engine changes. The key includes the data
        type, but only fp16 (Mamba-1, SSD) and bf16 (Gated DeltaNet) were
        measured: add one bf16-versus-fp16 SSD comparison (c, optional, RTX
        3090 lab host; under 1 GPU-hour) or state that the dtype dependence is
        untested. (R3 M4.5, m19; R2 m16, M4.3)
\end{enumerate}

\subsection*{M8. Clarity and self-containedness}
Merges EIC M5, R1 m12, m14, R2 M7, R3 M7, DA C12, RC M9a, M9c, ED reason 5.
The paper uses five or more timing conventions, packs 3--10 numbers into most
sentences, has two figures in 20 pages and depends on the supplement for the
evidence behind the recipe (Tables S33, S35, S36, Fig.~S3, the S2 transfer).

% STAGE: 10 (revision)
% RULES: academic-paper/agents/draft_writer_agent.md
\begin{enumerate}[label=M8\alph*.]
  \item (a) \textbf{Abstract.} About 180 words, three findings, one number per
        claim with its convention named (per call, latency-weighted, device
        time, graph replay); no undefined terms (``recurrence-preserving
        launch'', ``three-parameter ranker'', ``residual regret''); the
        in-house kernel disclosed (M2a). \emph{Where:} main.tex:47--70. (RC
        M9a; EIC M5; R2 M7; R3 M7.5; DA C12; R1 m12)
  \item (a) \textbf{Restructure Sec.~6 around the findings.} One subsection
        per question of Sec.~1 (or per finding of M1a), each with one table or
        figure and a one-sentence takeaway; split Table~5 into ``profiled
        selectors (information test)'' and ``calibration-only selectors'';
        move CMI, permutation tests, depth-6 trees and the bootstrap of the
        fit to the supplement; compress the citation lists of Sec.~2, the
        information tests of Sec.~6.3 and the key-churn derivation of
        Sec.~6.7 (its practical conclusion: profile at warm-up, 30--50\,ms per
        key). \emph{Where:} Sec.~6 (647--1395). (EIC M5; R3 M7.1, M7.4; R2 M7;
        R1 m14)
  \item (a) \textbf{One convention per table and a glossary.} State the
        convention in each column header and footnote exceptions: Table~4
        mixes device-time Mamba-1 rows with launch-inclusive Mamba-2 rows;
        Table~8 mixes device time, launch-inclusive and above-floor rows;
        Table~5 mixes sessions and seeds; Sec.~6.6 quotes launch-inclusive
        and device-time headroom in one sentence. Replace the Conventions box
        by a small table mapping each convention to the deployment regime it
        approximates and the tables that use it, and add a glossary (oracle
        headroom, headroom, residual regret, floor, covered, tie,
        latency-weighted share, form-blind, tuned shapes). \emph{Where:}
        main.tex:608--642; Tables 4, 5, 8. (EIC M5; R2 M7; R3 M7.2, m4; DA
        C12)
  \item (a)+(b) \textbf{Figures and main-text anchors.} Add a mechanism figure
        (latency versus chunk decomposed into the model's terms or the M5e
        profiler metrics, per regime) and a claims-to-evidence table; bring
        into the main text the cost-model fit (Fig.~S3) or the
        profiling-budget curve (Fig.~S4), Table S33 merged with Table~7, the
        S2 cross-device result, the S11 divergence result and the TP table
        (now prose only); every abstract and conclusion claim must point to a
        main-text table or figure, within V4. (R2 M7; R3 M7.3; EIC M5; RC
        M9c; ED)
  \item (a) \textbf{Discussion and Limitations.} Expand the Discussion
        (main.tex:1397--1412, two short paragraphs) to about one page:
        integration into Triton's autotuner and \texttt{do\_bench\_cudagraph},
        engine graph capture, generality beyond these kernels, and what to
        expect on data-centre GPUs; itemize Limitations (1416--1440, one
        run-on paragraph). (EIC M5)
\end{enumerate}

\subsection*{M9. Positioning and missing related work}
Merges R2 Sec.~7, R2 M3(a), R2 m9, DA C17, EIC M3, EIC m17. Verify the
bibliographic data of every added reference before citing it.

% STAGE: 10 (revision)
% RULES: academic-paper/agents/draft_writer_agent.md
\begin{enumerate}[label=M9\alph*.]
  \item (a) \textbf{Static and predictive pruning of GPU autotuning} (closest
        prior art for ``compile, read the resource report, rank, profile the
        top-$k$''): Ryoo et al., Program Optimization Space Pruning for a
        Multithreaded GPU (CGO 2008); Lim, Norris and Malony, Autotuning GPU
        Kernels via Static and Predictive Analysis (ICPP 2017); Bergstra,
        Pinto and Cox (InPar 2012); Falch and Elster (CCPE 2017); the survey
        of Balaprakash et al.\ (Proc.\ IEEE 2018); Kaufman et al.\ (MLSys
        2021). Occupancy and latency modelling for Eq.~(3): Volkov, Better
        Performance at Lower Occupancy (GTC 2010); Hong and Kim (ISCA 2010).
        \emph{Where:} Sec.~2 ``Performance models''; Sec.~4. (R2 M3(a), Sec.~7)
  \item (a) \textbf{Shape-keyed tuning in production stacks} (the recipe's
        closest practice): PyTorch~2 / TorchInductor (ASPLOS 2024;
        max-autotune, dynamic-shape bucketing and specialization, per-shape
        CUDA graphs); FlashDecoding++ (MLSys 2024; shape-keyed GEMM selection
        in serving); MIOpen find-db/perf-db (arXiv:1910.00078); cuDNN
        benchmark mode; TensorRT's persistent timing cache; the per-device,
        token-count-bucketed Triton configuration files of the fused-MoE
        kernels in vLLM and SGLang; SGLang (NeurIPS 2024); recent autotuners
        BaCO, GPTune and ytopt. Position the recipe against them and say what
        it adds. (R2 Sec.~7; DA C17)
  \item (a) \textbf{TFLA.} TFLA is called ``the closest analysis'' but appears
        only as a failing baseline on a scan it does not model, so its
        rejection of the quadratic term is foregone; its runtime form is
        essentially the paper's work-term model $BL(a+bC)+B\lceil
        L/C\rceil g$ plus a floor. State this, compare TFLA's expression with
        the work-term model on SSD and Gated DeltaNet, where it applies, move
        the TFLA-on-scan baseline to the supplement, and expand the two
        sentences at main.tex:221--231. (R2 M3(a), m9; EIC M3)
  \item (a) \textbf{Background.} Parallel scans for linear recurrences:
        Blelloch (1990); Martin and Cundy (ICLR 2018); Merrill and Garland
        (2016, decoupled look-back, the pattern of a chunked scan with a
        carried inter-tile state); Smith, Warrington and Linderman (ICLR
        2023). Measurement variability: Sinha et al., Not All GPUs Are Created
        Equal (SC 2022), for the per-device tuning and power-cap discussion.
        (R2 Sec.~7)
\end{enumerate}

% ======================================================================
\section*{Minor Issues}
\begin{enumerate}[label=N\arabic*.]
  \item Notation: main.tex:516 ``Tensor parallelism over $t$ devices runs the
        scan at width $d/n_{\mathrm{tp}}$''; use $n_{\mathrm{tp}}$ throughout.
        (EIC m4; R3 m1; RC M9b)
  \item Notation: Table S14 / appendix.tex:868 labels the first parameter
        $t_0$ where Eq.~(3) uses $\tau_0$, and $t_0$ is also the tile start in
        the pseudo-code. (R3 m2; RC M9b)
  \item Notation: $k$ denotes both a shape key (Eq.~(5), main.tex:1296--1299)
        and top-$k$; rename the key index (e.g.\ $\kappa$) or write top-$m$.
        (RC M9b)
  \item Notation: in Eq.~(3) waves $w$, resident warps $W$ and thresholds
        $W_0, W_1$ are easily confused; rename the waves (e.g.\
        $n_{\mathrm{wave}}$). (EIC m5)
  \item Caption length: Table~5 (103 words, main.tex:827--837), Table~6 (105
        words, 996--1005) and Table~3 (82 words); move definitions to the
        conventions table or text and keep captions under about 50 words.
        (RC M9e)
  \item Fig.~1: the caption omits the timing convention (launch-inclusive per
        Table S1, which explains the 0.5\,ms SSD and Gated DeltaNet floors and
        the 0.06\,ms Mamba-1 floor at $L=512$); the panel titles sit on the
        tick labels of the panel above; three stacked narrow panels leave half
        the width empty (use a $1\times3$ layout with a shared legend); the
        caption says chunk 512 is fastest at batch 1 while Table~3 has
        $C^*=256$ at $L=512$ (say ``tied'' or qualify by length). (R3 m6; EIC
        m7)
  \item Fig.~2: fonts too small for print, only three x-axis ticks; in the
        Triton panel ``default'' is a different chunk per curve; the host-bound
        RTX~3090 series is visible only in the Triton panel (say so if it lies
        at zero elsewhere) and a host-bound RTX PRO~6000 series is missing;
        annotate bins with call counts; add the latency-weighted series.
        (R3 m7; R2 m13)
  \item Fig.~S5: the rotated labels ``1/8 histogram'' and ``full histogram''
        overlap the plotted lines. (R3 m8)
  \item Table~2: the ``Calib.''\ column needs a unit (timings per GPU); split
        the ``Regret'' cell ``0.4*/3.5*/3.0 (4.2/6.2/3.0)'' into ``terms
        selected on these calls'' and ``a priori'' columns. (EIC m6; R2 m12;
        R3 m13)
  \item Table~4: write ``calls whose oracle differs from the best static
        chunk'' and ``entropy-rule regret'' instead of ``Oracle $\neq$
        static'' and ``Rule regret''. (R3 m12)
  \item Table~5: report MLP rows as five-seed mean $\pm$ s.d.\ (seed 0 gives
        0.37\% on the RTX PRO~6000, the five-seed mean 1.59\%); footnote the
        RTX~3090 $\dagger$ row's different session prominently or use the
        same session; the first row merges the default chunk 256 and the
        entropy rule, which coincide, and it coincides with ``Best single
        chunk, held out'' on two of three GPUs: give each selector its own
        row or explain the coincidence in one footnote; ``Best chunk per
        batch, held out'' is computed in sample in the stored analysis (the
        held-out values are identical, 20.0 / 17.4 / 24.7), so fix the label's
        provenance. (R1 m3, m15; R3 m9; EIC M2 fix (d))
  \item Table~7: give the oracle scan time (``(20.1)'') its own column; the
        work-term ``Eager / Refit'' columns use per-call regret, unlike their
        neighbours: separate them visually or move them to the supplement.
        (R3 m10; EIC m15)
  \item Table~8: rows mix conventions (see M8c); add the oracle scan time per
        forward pass, since the identical 4.42\,ms default regret on the RTX
        PRO~6000 and the H800 is correct (oracle 13.1 and 24.2\,ms) but looks
        like a copy error. (R3 m11)
  \item Reconcile one quantity across tables: RTX PRO~6000 device-time cost
        model 0.40\% / p95 3.8 in Table~5, Sec.~6.4 and Tables S30, S31
        (0.405 on round medians) against 0.38\% / 3.6 in Tables S14, S15, S16,
        S19 (0.383 on round means); MLP 0.37 vs.\ 0.35; top-2/top-3 0.23/0.04
        vs.\ 0.21/0.02 (RTX~3090 0.72/0.55 vs.\ 0.67/0.51); RTX~3090 cost model
        3.52 vs.\ 3.55 (two sessions) and three-parameter ranker 0.85 vs.\
        0.88. Use one oracle definition or state which each table uses. (EIC
        m3; R1 m1; DA C13)
  \item The RTX~3090 ordinary least-squares regret is 1.74\%, printed as
        1.8\% (appendix.tex:847, and main.tex wherever repeated); write 1.7\%.
        (RC M3c)
  \item main.tex:714: ``In 94\% of the 414 cells the input-agnostic chunk is
        tied'' uses the CI-overlap criterion (94.4\%); under the 2\% rule the
        Conventions box makes default it is 93.5\%. State the criterion or
        report 93\%. (EIC m1)
  \item main.tex:814--815 mixes populations: the held-out 0.14\% is over all
        2360 calls and depends on the fallback rule for the 204 uncovered calls
        (a plain length-bucket fallback gives 0.22\%); 0.14\,$\mu$s is over
        covered calls only (0.37\,$\mu$s over all calls). Specify the
        population and the fallback rule. (EIC m2)
  \item Reference-scan stock dispatch: 6.6\% (Table~6) compares a separately
        timed stock binary against an oracle from the runtime-chunk build; the
        runtime-chunk build at the stock chunk gives 5.35\%, same-chunk
        re-timings across the two binaries differ by 0.86--1.13 (10th--90th
        percentile), and the device-time stock pairing is 3.1\% (Table S21).
        Use the same binary and repeated rounds, and say which number supports
        ``a shape tree halves that''. (R1 m2; EIC m14)
  \item Abstract ``three GPU generations'': only the Triton Mamba-1 scan was
        measured on Hopper; the production-kernel results cover two GPUs. (EIC
        m11; R3 m14)
  \item ``Workstation GPUs'': the RTX PRO~6000 is the Blackwell Server
        Edition (Table S2) and the RTX~3090 a consumer GeForce board; write
        ``an Ampere consumer GPU and a Blackwell server GPU''. (R2 m2; DA C14;
        RC M6a)
  \item Abstract qualifier ``0.4 to 3.5\% when its terms are selected on the
        same GPUs' calls'' includes the H800 (3.0\%), whose calls did not
        select the terms. (DA C14)
  \item Gated DeltaNet is limited by the library's solve\_tril path to chunks
        $\le64$ (a three-point design space); say in the abstract, Sec.~6.5
        and the conclusion that its ``fixed optimum at 64'' may be a library
        ceiling (the work-term model places the optimum at 61--85). (EIC m11;
        R2 m7)
  \item The ``compile-facts rule (no parameters)'' of Table S (topk\_all) is
        never defined; define it. (RC M5a)
  \item ``Every configuration produced the greedy tokens of the unmodified
        model'' (main.tex:1194) refers to three prompts, while S11 shows
        2.6--4.7\% of 192 prompts diverging for fixed alternative Mamba-1
        chunks in fp16; add ``on these three prompts''. (RC further
        observations)
  \item The ``shape'' feature set of Sec.~6.3 includes the layer index, which
        is not part of the key of Sec.~3; the regrets are identical without it
        (0.465 / 0.383 / 1.626 / 0.732), so describe the feature set
        accurately or drop the layer. (DA C15)
  \item Supplement citations ``Table S33'' and ``Section S11''
        (main.tex:1409, 1412) lack the ``Supplementary'' prefix used
        elsewhere. (EIC m10; R3 m17)
  \item Sec.~1, third question: ``\ldots\ the trade-off between model-driven
        and empirical optimization [Yotov et al.\ 2003]?'' is ungrammatical;
        rephrase. (R3 m18)
  \item Reference-scan dispatch: ``selects its chunk \ldots\ from the length
        at compile time'' is inaccurate; the stock scan selects among
        compile-time template instantiations at run time (host-side dispatch
        on the sequence length). Align with Table~1's caption. (R2 m1)
  \item Chunk-spread statistics (8.5$\times$; 15.2$\times$ at p95) are
        dominated by heavily spilling configurations (e.g.\ $D_b=16$ with two
        warps); demote them behind the regret of the default or best static
        choice. (R2 m4)
  \item ``Compile facts'' is idiosyncratic; use ``compiler resource report''
        or ``static resource usage''. (R2 m17)
  \item The 2\% tie criterion is below launch-inclusive noise (CV 4--5\%);
        reference the three-criterion tie table (Table S3) wherever
        launch-inclusive tie shares appear. (R1 m10)
  \item Add the latency-weighted share of the batch-1 static chunk per GPU
        (15.1 / 6.2 / 7.2\%) to the supplementary static/bucket table; Table~4
        gives only the share of chunk 512. (RC M4b)
  \item Bibliography: the cited preprint is filed under cs.CV (check the
        entry); update 2026 arXiv-only references (Mamba-3, log-linear memory,
        Astra) to archival versions where they exist; fill empty
        publisher/address fields; add pages for beck2025tfla and
        dao2024flashattention2; verify the page range 73099--73134 of
        wang2026tilelang. (EIC m17; RC Stage 11 gate)
  \item Typography: set the Conventions and Recipe boxes as lightly shaded
        boxes in normal-size text instead of \texttt{\textbackslash
        footnotesize} inside a tight \texttt{\textbackslash fbox}; Tables 5
        and 7 use \texttt{\textbackslash scriptsize} with 2--2.2\,pt column
        separation. (R3 m24, M7)
\end{enumerate}

% ======================================================================
\section*{Acceptance Criteria}
The paper moves to Minor Revision when A1--A8 hold; to Accept when, in
addition, M1d is settled by measured end-to-end evidence (M3b at minimum) and
every remaining claim is scoped to the evidence.
\begin{enumerate}[label=A\arabic*.]
  \item \textbf{Compliance (V1--V5).} A clean, current anonymized artifact
        backs every main and supplementary table (an anonymization scan finds
        no host names, user paths, hashes, project or venue names); only the
        two PDFs are uploaded; the preprint relationship is disclosed to the
        editor; the main text stays within 20 pages excluding references and
        the supplement is about 20--25 pages.
  \item \textbf{Framing (M1).} The contributions are the robust findings of
        M1a; the data-dependence refutation is one subsection with the lemma
        as an observation; the headroom result is stated as an information
        result with a host-side cost corollary, including the RTX PRO~6000
        reduction cost; the significance claim is settled per M1d.
  \item \textbf{Kernel and baselines (M2).} The in-house kernel is disclosed
        as a stress case in the abstract and Introduction; the free length
        rule, the best chunk per batch and a length-bucket table are the
        untuned references in the abstract, Sec.~6 and Tables 4, 5, 8; one
        definition of default and best static per kernel; the chunk effect at
        fixed launch configuration is reported and its slice re-timed under
        the full protocol.
  \item \textbf{Device-bound evidence (M3).} Recoverable gains with cluster
        intervals replace oracle-relative shares; a full-model prefill
        captured as one CUDA graph (default vs.\ table, two sessions) gives the
        measured end-to-end effect and scan share; packed steps are re-timed
        under graph replay with a noise floor and the recipe and
        Contribution~2 follow the result; engine execution modes are described
        with sources; the engine-level and data-centre items are either run or
        their absence is stated and the claims scoped to the GPUs measured.
  \item \textbf{Cost model (M4).} One pre-specified a-priori protocol with all
        six reported, latency-weighted regret, maxima, per-batch values and
        intervals; the abstract and conclusion state the protocol-dependent
        accuracy with its latency-weighted companion; the form-blind
        comparison with the three-parameter restriction and the MLPs is stated
        plainly; the fit is documented; the latency term is justified or
        refitted; the model is positioned as a cross-device or pruning case
        study.
  \item \textbf{Measurement and statistics (M5).} Noise floors next to every
        regret; the RTX~3090 TP numbers demoted or re-measured; SSD and Gated
        DeltaNet value invariance measured above the floor (or the recipe
        qualified); H800 input effects in the main text; counters support the
        stated mechanism where the profiler is permitted; the stock lab-host
        corroboration reported; every main-table number from a guarded run or
        flagged; one value per quantity across all tables.
  \item \textbf{Workload and recipe (M6, M7).} Workload-weighted headline
        numbers beside unweighted ones; the serving replay costed with the
        production kernels' measured timings, with measured and modelled
        conclusions labelled; the routed study led by stock dispatch and
        prefill; a recipe that handles batch invariance, excludes chunks that
        overflow 32-bit indexing, keys on the specialization signature with an
        alignment/padding action, qualifies synthetic warm-up inputs, and
        labels untested items as precautions.
  \item \textbf{Clarity and positioning (M8, M9, N1--N34).} An abstract of
        about 180 words with one labelled number per claim; Sec.~6 organized
        by finding with a mechanism figure and a claims-to-evidence table; one
        convention per table and a glossary; an expanded Discussion and
        itemized Limitations; the missing prior art cited, verified and
        positioned against; all minor issues resolved.
\end{enumerate}

\paragraph{Estimated new GPU time.} Minimum for A3, A4 and A6 on the RTX~3090
lab host: about 20--30 GPU-hours (M2d re-timing 2--3, M3b 4--8, M3d 3--5, M4f
1, M5c 3--4, M5e 2--4, M5f 3--5, M5i under 1, M7d/M7e under 1 each), i.e.\
about 5--8 hours of wall time on its four GPUs. Optional lab-host items (M2f,
M3g, M4g, M5b re-run, M6b, M7f dtype): about 15--25 GPU-hours more. Optional
items that need a rented GPU (M3e engine integration 15--25, M3f data-centre
replication 6--10, M4h fourth GPU 2--4, M5d H800 control 1--2, and a second RTX
PRO~6000 session of about one GPU day): about 30--60 GPU-hours; not planned in
this cycle by the author's decision. All other items are text or offline
re-analysis of existing result files.

\end{document}
